\section{Motivation \& Background}

We want to : 1) run programs that can interact with the outside world.
2) formalize the guarantees that we expect the hardware to provide
3) have good performance.

\subsection{Security Monitor}
\label{sec:securitymonitor}
In many enclave platforms, including Citadel, the Security Monitor (SM) \cite{costan2016sanctum} \cite{lebedev2019sanctorum} is a small (9K LOC), trusted piece of software running at a higher privilege mode than the hypervisor or the OS.
Its role is to link low-level invariants exposed by the hardware (e.g., if a specific memory access is authorized), invoking the right microarchitectural cleanup routine, and the high-level security policies defined by the platform (e.g., which enclave is currently executing on which core and which memory region does it own).

\subsection{Microarchitectural Timing Side Channels}
Microarchitectural side channels allow for information to be leaked from one security domain to another.
These attacks exploit microarchitectural structures shared between different security domains, such as memory caches \cite{gotzfried2017cache, brasser2017software}, translation look-aside buffers (TLBs) \cite{gras2018translation}, branch predictors \cite{evtyushkin2018branchscope, aciiccmez2006predicting}, DRAM controllers \cite{wang2014timing}, etc., to transmit a signal.
Because microarchitectural state cannot be directly observed through the Instruction Set Architecture (ISA), an attacker will typically rely on the timing of events to indirectly observe the microarchitectural state.
For instance, if a piece of data has been loaded in the cache, accessing it will be faster than if it is only stored in main memory.
Using various methods to prepare the shared microarchitectural state and observe the state changes due to the victim's activity, the attacker can recover sensitive information  about the victim's execution, for example, which addresses have been accessed, which branches were taken in which directions, and/or information about victim program's inputs.
Cache side channels are the easiest channels to exploit, and so have been extensively studied, but any piece of shared microarchitectural state can theoretically be used as well.

To summarize the observations that could be made through side-channels, the community formalizes side-channels as a set of observations that an attacker could make about the microarchitectural state. 
For example, cache side-channels on unprotected machines typically lead to an observation model that includes the sequence of addresses of memory operations (sometimes with their timestamps) as part of the observation.

\subsection{Transient Execution Attacks and Execution Models}
\label{sec:specandside}

Speculative or transient execution is one of the key performance features of modern processors.
Unfortunately, it can make reasoning about security boundaries more difficult.
Some microarchitectural data structures that orchestrate speculation (e.g., branch predictors) can be trained so the processor will execute code snippets speculatively, in ways that are incorrect or on data that should not be accessed during normal execution.
Concretely, instructions that were not part of the program normal execution path can be executed out-of-thin-air, speculatively, by the processor.
An example is the execution of code speculatively before the resolution of a preceding security check, such as a memory access bound.
These erroneous executions happen speculatively, and these instructions will be squashed by the processor.
Nonetheless, they will still change the microarchitectural state of the machine.
For example, a speculative load to some data will not affect the values of the architectural registers but will load the data in the memory hierarchy and possibly change the state of caches.
This means that a program that manages sensitive data can be coerced into leaking private data speculatively through a side channel.
More importantly, once some specific code snippet, or gadget, has been identified, these attacks can be used to leak virtually any data located in the victim's virtual address space.
This type of attack, called Spectre, was first introduced in \cite{Kocher2018spectre},  and many variants have since been demonstrated.